% method_mpc.tex
% Methods: the model-predictive rebalancing controller and the backtest design.
% Sources (the code is the truth where docs disagree): code/src/StudyBacktest.jl (run_study, _study_problem,
% _study_decision), code/src/Triggers.jl (check_trigger), code/src/Projection.jl (project_gbm),
% code/src/Covariance.jl (growth_covariance, shrinkage_loss), code/src/OnlineSIM.jl, code/src/SIM.jl,
% code/src/MarketInputs.jl, code/src/Screens.jl, code/src/Bandit.jl (_stream_seed), code/src/StudyRunner.jl
% (run lists, study_stats), code/src/GrowthScorecard.jl, code/src/StudyTables.jl (passive_references),
% config/study.toml, results/calibration_summary.toml, README.md, scripts/dev/gamma_seeds.jl,
% scripts/dev/large_k.jl, scripts/30_talk_data.jl, code/src/TalkData.jl.

\subsection{Model-predictive rebalancing}\label{sec:method-mpc}

The controller decides when to re-plan the book and leaves it alone in between. It follows model-predictive
control in the receding-horizon sense~\cite{garcia1989mpc,mayneEtAl2000}. At each re-plan it measures the
state of the account (the shares $n_i$, the cash $c$, the wealth $W(t)$, the close prices $S_i(t)$ and the online
estimates of that day), solves the constrained utility maximization of Algorithm~\ref{alg:allocation} for target
% config/study.toml [triggers] T = 21
shares $n^\star$, and predicts the distribution of the wealth of holding $n^\star$ over the next $T = 21$ trading
days with a multi-asset geometric Brownian motion (GBM). It then compares the measured wealth with the prediction
every day and re-plans from the newly measured state when the horizon elapses, when the wealth leaves the
predicted band, or when the book falls a set fraction below its peak. The plan over the horizon is to hold
$n^\star$ without further trades, and the GBM is the controller's only model of prices. The plant is the price
process that the plan meets, the recorded prices in the backtest and synthetic prices from per-asset JumpHMM
generators~\cite{alswaidan2026variance} in the stress test, where the controller still forecasts with the GBM and
never sees the generator. Multi-period trading by convex optimization~\cite{boydEtAl2017} also re-solves from
the measured state, but it plans a sequence of trades over the horizon and executes the first of them every
period. Our controller instead solves a single-period allocation, uses the forecast as a reference for the
realized wealth, and trades only when a trigger fires, so the band around the forecast acts as a no-trade
region~\cite{davisNorman1990} placed on the portfolio value relative to its forecast rather than on the weights.

At a re-plan on day $t$ the drift of name $i$ is its expected growth rate
$\mu_{g,i} = \alpha_i(t) + \beta_i(t)\,\tilde g_M(t)$ ($\mathrm{yr}^{-1}$). Here $\alpha_i(t)$ ($\mathrm{yr}^{-1}$)
and $\beta_i(t)$ (dimensionless) are the online single-index estimates of name $i$ against SPY~\cite{sharpe1963},
% config/study.toml [sim] ewls_half_life = 63, ewls_prior_weight = 63; [market] gtilde_window = 252
fitted by exponentially weighted least squares with a half-life of $63$ trading days from a prior worth $63$
pseudo-observations at the calibration fit (Appendix~\ref{app:controller}), and $\tilde g_M(t)$
($\mathrm{yr}^{-1}$) is the exponential moving average of the daily growth rate of SPY with a span of $252$ trading
days. The covariance $\Sigma_g$ ($\mathrm{yr}^{-2}$) of the daily growth rates of the names in the allocation
problem is
\begin{equation}
  \Sigma_g \;=\; \sigma_M^2(t)\,\beta\beta^{\top} \;+\; (1-\lambda)\,\hat\Omega_\varepsilon
  \;+\; \lambda\,\mathrm{diag}\big(\hat\Omega_\varepsilon\big),
  \label{eq:mpc-cov}
\end{equation}
where $\beta$ is the vector of the loadings $\beta_i(t)$, $\sigma_M(t)$ ($\mathrm{yr}^{-1}$) is the exponentially
weighted standard deviation of the daily growth rate of SPY with the same $63$-day half-life, and
% config/study.toml [covariance] omega_window = 126
$\hat\Omega_\varepsilon$ ($\mathrm{yr}^{-2}$) is the sample covariance of the residuals
$\varepsilon_i(s) = g_i(s) - \alpha_i(t) - \beta_i(t)\,g_M(s)$ over the $L_\Omega = 126$ daily growth rates that
end on day $t$. The weight $\lambda \in [0, 1]$ shrinks the residual covariance toward its diagonal, from the full
sample covariance at $\lambda = 0$ to the covariance of the single-index model, the target of Ledoit and
Wolf~\cite{ledoitWolf2003}, at $\lambda = 1$. We chose it on the calibration window by the QLIKE
% results/calibration_summary.toml (lambda = 0.0)
loss~\cite{patton2011} of $21$-day variance forecasts for random portfolios, which selected $\lambda = 0$
(Appendix~\ref{app:controller}). The allocation problem uses the same $\mu_g$ and $\Sigma_g$ in its volatility limit
and in its minimum-variance and tangency benchmarks, so each book is forecast with the model that sized it.

% config/study.toml [triggers] n_paths = 1000, T = 21
From the target shares $n^\star$ and the close prices $S_i(t)$, the controller simulates $N_p = 1{,}000$ paths over
the next $T = 21$ trading days. On path $p$ and day $\tau = 1, \dots, T$ after the re-plan, the log price of each
name advances by $\Delta t$ times a vector of growth rates drawn from a multivariate normal distribution with mean
$\mu_g$ and covariance $\Sigma_g$, independently across days and paths,
\begin{equation}
  \ln S^{(p)}_i(\tau) \;=\; \ln S^{(p)}_i(\tau-1) \;+\; g^{(p)}_i(\tau)\,\Delta t,
  \qquad g^{(p)}(\tau) \sim \mathcal{N}\big(\mu_g, \Sigma_g\big),
  \qquad S^{(p)}_i(0) = S_i(t),
  \label{eq:mpc-gbm}
\end{equation}
a multi-asset GBM~\cite{glasserman2003} written in growth rates. Because $\mu_g$ is the expected growth rate of the
log price, the drift needs no It\^o correction. The wealth of the plan on path $p$ is
\begin{equation}
  W^{(p)}(\tau) \;=\; \sum_i n^\star_i\, S^{(p)}_i(\tau) \;+\; c^\star\, e^{g_f(t)\,\tau\,\Delta t},
  \qquad c^\star \;=\; W(t) - \sum_i n^\star_i\, S_i(t),
  \label{eq:mpc-value}
\end{equation}
where $c^\star$ (USD) is the cash left after paying for the targets at the close of the decision day. The forecast
is the sample mean $\mu_W(\tau)$ and standard deviation $\sigma_W(\tau)$ (USD) of $W^{(p)}(\tau)$ across the paths.
It omits the trading costs, the overnight price change before the fill and, for a 130/30 book, the borrow fee, and
its paths are seeded with the run seed and the decision date, so a run reproduces its forecasts exactly.

% config/study.toml [triggers] z = 1.96, D_max = 0.08, drawdown_reset = true
The forecast defines a band $\mu_W(\tau) \pm z\,\sigma_W(\tau)$ with $z = 1.96$, which would contain the wealth on
about $95\%$ of days if the forecast wealth were normal and the model were correct. With $\tau$ the number of trading
days since the last re-plan, the controller marks the wealth $W(t)$ at each close, after that day's fills, costs,
cash growth and fees, updates its running peak $W^{\mathrm{pk}}(t)$, and checks three triggers in turn, acting on
the first that holds. The drawdown trigger fires when $1 - W(t)/W^{\mathrm{pk}}(t) > D_{\max}$ with $D_{\max} = 0.08$, the horizon
trigger when $\tau \ge T$, and the band trigger when $W(t)$ lies outside $\mu_W(\tau) \pm z\,\sigma_W(\tau)$, so the
wealth measured $\tau$ days after a re-plan is compared with the forecast for $\tau$ days ahead on days
$\tau = 1, \dots, T-1$. A drawdown re-plan sets the peak back to the current wealth, so the trigger fires once for
each fresh loss of $D_{\max}$ rather than on every day that the book stays below an old peak.
% README.md (88.5% of Tier 1 decisions before the reset of 2026-10-25)
We adopted this reset after an earlier pass of the Tier 1 grid (the three per-sector selection rules crossed with
the seven weighting rules) without it, in which $88.5\%$ of the decisions were
% code/src/StudyBacktest.jl (_study_lots, trig.projection = nothing); config/study.toml cash_revisit = 21
drawdown re-plans. A re-plan that leaves no position worth more than $0.1\%$ of wealth sets no band, and the next
re-plan then comes from the drawdown trigger or, after $21$ trading days, from a cash revisit.

% config/study.toml [screens] repick_interval = 21, repick_on_band_exit = true
The selection rule picks the basket on the first day, every $21$ trading days after its previous pick, and at
every band exit. With the hedge on, each pick also refreshes the hedge pool, and for the 130/30 benchmark the short
candidates. When a timer pick changes the set of names on a day when no trigger fired, the controller re-plans as
well. Because band exits depend on each run's own wealth path, the dates of the picks, and with them the
baskets, could differ across weighting rules within a seed. At a re-plan it solves the allocation problem from the shares actually held, so the turnover limit counts
the changes from that book, queues orders to $n^\star$ for the next open, forecasts the new book from the measured
wealth $W(t)$ and resets $\tau$ to $0$. Each forecast therefore starts from the measured wealth rather than from
the previous forecast, and the horizon recedes with every re-plan. If the solve fails, the controller keeps every
holding, queues no orders and forecasts the kept book. One episode of a real run, three consecutive forecasts and
the re-plans that ended them, shows the mechanism (Fig.~\ref{fig:mpc}), and Algorithm~\ref{alg:mpc} gives the daily
loop of a run.

\begin{algorithm}[t]
\caption{Daily loop of one backtest run. A run fixes a selection rule, a weighting rule, the hedge setting and a
seed. Every decision uses data dated on or before the close of its day.}
\label{alg:mpc}
\begin{algorithmic}[1]
\REQUIRE trading days $t_1, \dots, t_{420}$; calibration fit on 2014--2024; $W_0$; $T$, $z$, $D_{\max}$, $N_p$
\ENSURE daily wealth $W(t)$ and the log of re-plans
\STATE $c \gets W_0$, $n \gets 0$, $W^{\mathrm{pk}} \gets W_0$, $\tau \gets 0$; no band, no pending orders
\FOR{each trading day $t = t_1, \dots, t_{420}$}
  \IF{orders are pending}
    \STATE fill them at the open: sells first, then buys scaled by one factor in $[0,1]$ so that $c \ge 0$; pay
      the costs of Eq.~\eqref{eq:cost}
  \ENDIF
  \IF{$t > t_1$}
    \STATE $c \gets c\,e^{g_f(t)\Delta t}$ if $c > 0$; charge the borrow fee on short positions; $\tau \gets \tau + 1$
  \ENDIF
  \STATE mark $W(t) \gets c + \sum_i n_i S_i(t)$ at the close; $W^{\mathrm{pk}} \gets \max(W^{\mathrm{pk}}, W(t))$
  \STATE update $\alpha_i, \beta_i, \sigma_M, \tilde g_M, \xi$ and $\gamma_i$ with the growth rates into day $t$
  \STATE reason $\gets$ the first that holds of: first day; cash revisit (last re-plan went to cash and
    $\tau \ge 21$); drawdown ($1 - W(t)/W^{\mathrm{pk}} > D_{\max}$); horizon ($\tau \ge T$); band exit
    ($|W(t) - \mu_W(\tau)| > z\,\sigma_W(\tau)$); otherwise none
  \IF{first day, or $21$ trading days since the last pick, or reason is band exit (except on the fixed schedule)}
    \STATE pick the basket with the selection rule; refresh the hedge pool and the short candidates
    \IF{reason is none and the set of names changed}
      \STATE reason $\gets$ basket change
    \ENDIF
  \ENDIF
  \IF{reason is not none}
    \STATE build $\mu_g$ and $\Sigma_g$ (Eq.~\eqref{eq:mpc-cov}) for the names of the allocation problem
    \STATE solve the allocation problem (Algorithm~\ref{alg:allocation}) from $n$, $W(t)$, $S(t)$: targets $n^\star$
    \IF{the solve failed}
      \STATE $n^\star \gets n$; queue no orders
    \ELSIF{$t < t_{420}$}
      \STATE queue orders to $n^\star$ for the open of the next day, each target worth at most $0.001\,W(t)$ set
        to zero shares
    \ENDIF
    \STATE simulate $N_p$ paths of Eqs.~\eqref{eq:mpc-gbm} and~\eqref{eq:mpc-value} for the book $n^\star$; store
      $\mu_W(\tau)$, $\sigma_W(\tau)$, $\tau = 1, \dots, T$
    \STATE set the band from the forecast (no band if the book holds no position above $0.001\,W(t)$);
      $\tau \gets 0$
    \IF{reason is drawdown}
      \STATE $W^{\mathrm{pk}} \gets W(t)$
    \ENDIF
  \ENDIF
\ENDFOR
\end{algorithmic}
\end{algorithm}

\subsection{Backtest design}\label{sec:method-backtest}

A backtest run fixes a selection rule, a weighting rule, the hedge setting and a seed, and replays
% config/study.toml [grid] n_seeds = 30, base_seed = 5660
Algorithm~\ref{alg:mpc} over the $420$ trading days of the trading window. We used the seeds $5660$ to $5689$, $30$
runs per cell of an experiment grid. The seed enters in three places: the random selection rules draw each basket
from a stream seeded with the run seed and the pick date, each sector's bandit seeds its exploration with the run
seed, the pick date and the sector, and the forecast paths of every re-plan are seeded with the run seed and the
decision date. For the deterministic selection rules (the $\gamma$-based rules and the hand-picked lists) the seed
changes only the forecast paths, and through them the band exits and so the dates of the re-plans and picks. Each
run with the hedge on has a twin with the hedge off and the same seed, identical except that the beta cap and the
hedge pool are removed.

% results/*_summary.csv and results/calendar_baseline/calendar_baseline_summary.csv: every row status ok; distinct
% labels over the eight files = 18,238 (14,038 + 4,200 calendar-only)
We ran nine grids over the full trading window, and all $18{,}238$ distinct runs finished
% results/tier1_summary.csv (854 rows)
(Table~\ref{tab:S-grids}). The Tier 1 grid crossed the seven weighting rules with three per-sector selection rules
% results/fixed_summary.csv (427 rows)
($3$ names per sector, $33$ names) and the hedge on and off ($854$ runs). The fixed-schedule grid repeated its
hedge-on half with basket picks only on the $21$-day timer ($427$ runs), while band exits and drawdown stops still
triggered re-plans, unlike the calendar-only grid. Because the random and bandit streams depend
only on the seed and the pick date, and the $\gamma$ ranking and the hedge pool only on market data, every weighting
rule in that grid held the same baskets within a seed.
% results/universe_summary.csv (4,438 rows = 4,410 universe-wide + 28 hand-picked)
The universe-wide grid applied five selection rules that pick $K$ names from the whole universe, at
$K = 5, 8, 13, 21$ and $34$ ($4{,}410$ runs), and the hand-picked grid ran the two hand-picked lists ($28$ runs).
% results/gamma_seeds/gamma_seeds_summary.csv (4,200 rows; 140 labels repeat tier1, universe-wide and hand-picked runs: 14 + 70 + 28 + 28)
% results/large_k/large_k_summary.csv (1,680 rows)
A seed rerun gave each deterministic rule all $30$ seeds ($4{,}200$ runs, $140$ of which repeat earlier runs), and
the large-book grid extended global $\gamma$ and random picks from the whole universe to $K = 100$ and $200$
($1{,}680$ runs).
% results/calendar_baseline/calendar_baseline_summary.csv (4,200 rows, z = Inf, D_max = Inf; n_horizon = 19 and
% n_basket_change = 0 in every run); scripts/dev/calendar_baseline.jl
A calendar-only grid repeated the $4{,}200$ runs of the main comparison, described next, with the band exit and the
drawdown stop switched off ($z = \infty$ and $D_{\max} = \infty$), so that a run re-planned only when the $21$-day
horizon elapsed, and paired each run with the triggered run of the same selection rule, weighting rule, hedge
setting and seed ($4{,}200$ runs).
% results/eta_summary.csv (2,562 rows); results/limits_summary.csv (17 rows)
Two sensitivity grids varied the CES elasticity $\eta$, one on the Tier 1 grid with the single-name cap on and off
($2{,}562$ runs) and one on the hand-picked $13$-name list under the budget alone ($17$ runs). This paper does not
report results from the fixed-schedule grid or from the elasticity grid on the Tier 1 rules, and the budget-only grid
served as the solver check against the closed forms.

The main comparison crosses ten selection rules with the seven weighting rules, with the hedge on and off. The ten
are $\gamma$, bandit and random per sector ($33$ names each), global $\gamma$, correlation-aware $\gamma$, liquid
then $\gamma$, random from the $100$ most liquid names and random from the whole universe ($K = 13$ each), and the
two hand-picked lists. The $13$-name list (AAPL, MSFT, INTC, MU, AMD, NVDA, GS, BAC, WFC, C, F, GM, JNJ) and its
$15$-name extension with WMT and PG were chosen in advance by the author, from sectors the author follows and
without 2025 data. Each is one list chosen by a person. Each of the $140$ cells holds $30$ runs, and we report the
median over the runs of each cell, never a median pooled across selection rules. For each run we computed the
terminal wealth ratio $W_T/W_0$, with $W_T$ the wealth at the close of 2026-09-04, the realized growth rate
$\ln(W_T/W_0)/(419\,\Delta t)$, the maximum drawdown $\max_t\big(1 - W(t)/\max_{s \le t} W(s)\big)$, and the
annualized Sharpe ratio~\cite{sharpe1966} $\bar y/(\sqrt{\Delta t}\,s_y)$, where $\bar y$ and $s_y$ are the mean and
standard deviation of the daily excess growth $y(t) = g_W(t) - g_f(t)$, with $g_W$ the daily growth rate of the
run's wealth. We also recorded the intercept $\alpha_W$ ($\mathrm{yr}^{-1}$) and slope $\beta_W$ of an ordinary least-squares regression of $y(t)$ on the daily excess growth
of SPY, and the number of re-plans by trigger. Two passive benchmarks received the same metrics. SPY buy-and-hold
puts $W_0$ into SPY at the close of 2025-01-02, and the equal-weight universe puts $W_0/373$ into each name at that
close and never rebalances. Both are valued at the closes and pay no trading costs. The universe-wide selection
rules, their thresholds and the grid of $K$ were designed after the Tier 1 results were known, so their comparison
is exploratory. The $30$-seed reruns and the large-book runs were added afterwards to measure how much the forecast
seed moves a deterministic rule and whether the comparison of global $\gamma$ with random picks holds for larger
books. The calendar-only grids, on real prices and on the synthetic stress sets, were added last, after all other
results were known, to measure what the two early triggers changed.

